\section{A factorization-free squarefree producer for nonzero norms}
\label{sec:squarefree-norm-audit}

The following audit establishes a useful automatic nonvanishing producer.
Its coefficient algebra has zero divisors in general; reducedness, rather
than an unjustified field assumption, is the property needed by the proof.

\begin{theorem}[Squarefree input supplies the norm witness]
\label{thm:squarefree-norm-producer}
Let \(K\) be a field, let \(f\in K[X]\) be nonconstant with
\(\gcd(f,f')=1\),
and let \(m\geq1\).  For every nonzero reduced polynomial
\[
 G(X,Y)=\sum_{j=0}^{m-1}A_j(X)Y^j,
\]
the multiplication norm in \(K[X,Y]/(Y^m-f(X))\) is nonzero.  Moreover,
writing \(d=\deg f\),
\[
 \deg N(G)\leq
 W:=\max_{A_j\ne0}\{m\deg A_j+jd\}.
\]
No hypothesis \(\gcd(m,d)=1\), no factorization of \(f\), and no
\(K\)-rational root of \(f\) is required.  The result holds in every
characteristic, including characteristics dividing \(m\).
\end{theorem}

\subsection{Reducedness with explicit divisibility witnesses}

Here the effective squarefree condition means the separability certificate
\begin{equation}
 \label{eq:squarefree-bezout}
       U f+V f'=1.
\end{equation}
For polynomial coefficients in an exact field, the ordinary degree-bounded
Euclidean algorithm computes this certificate when \(\gcd(f,f')=1\).
Over \(\mathbb Q\), or any perfect field, this is equivalent to the usual
absence of repeated irreducible factors.  Over an imperfect field of
positive characteristic those formulations differ: \(X^p-t\) over
\(\mathbb F_p(t)\) is irreducible but has zero derivative.  The present
simple-root argument uses the displayed Bezout certificate, not an
unqualified ring-theoretic \texttt{Squarefree} assumption.

\begin{lemma}[A factorization-free reducedness proof]
\label{lem:quotient-reduced-explicit}
For \(n\geq1\), if \(f\mid b^n\), then \(f\mid b\).
Consequently \(K[X]/(f)\) is reduced.
\end{lemma}

\begin{proof}
Compute the monic gcd \(g=\gcd(f,b)\), the quotients
\(f=gh\) and \(b=gq\), and a Bezout certificate
\(Af+Bb=g\).  Since \(g\ne0\), cancellation in the polynomial domain
gives
\begin{equation}
 \label{eq:hq-bezout}
                   Ah+Bq=1.
\end{equation}
Differentiate \(f=gh\) and substitute in
\eqref{eq:squarefree-bezout}.  This gives the explicit identity
\[
 1=h(Ug+Vg')+g(Vh').
\]
Put \(C=Ug+Vg'\) and \(D=Vh'\); hence \(Ch+Dg=1\).
Multiplying this identity by \eqref{eq:hq-bezout} yields
\begin{equation}
 \label{eq:hb-bezout}
  1=Eh+Fb,
  \quad E=CAh+CBq+DAg,
  \quad F=DB.
\end{equation}
Indeed the remaining product term is \(DBgq=Fb\).

The hypothesis implies \(b^n=hr\) for an explicit quotient \(r\).
Using \(1-Fb=Eh\), the finite geometric identity gives
\[
 \begin{split}
 1&=(Fb)^n+(1-Fb)\sum_{i=0}^{n-1}(Fb)^i\\
  &=h\left(F^n r+E\sum_{i=0}^{n-1}(Fb)^i\right).
 \end{split}
Write the displayed parenthesis as \(J\).  Then \(hJ=1\), and
\[
                  b=gq=(gh)(qJ)=f(qJ).
\]
This is a concrete divisibility witness, not merely an appeal to prime
factorization or unique factorization.  The only variable-length identity
is a finite geometric sum, proved by quantifier-free induction on \(n\).
No division by the numeral \(n\) is used, so the argument is valid in
positive characteristic.
\end{proof}

\subsection{A formal simple root in a computable coefficient ring}

Let
\[
 R=K[U]/(f(U)),\qquad u=U\bmod f.
\]
For implementation, reduce residues modulo the monic polynomial
\(\widehat f=f/\operatorname{lc}(f)\).  This changes the division
algorithm only: the original \(f\), rather than \(\widehat f\), remains
the polynomial in the norm equation \(Y^m=f(X)\).

The algebra \(R\) is nonzero because \(\deg f\geq1\), and the natural
map \(K\to R\) is injective.  In particular a nonzero scalar coefficient
remains nonzero.  By Lemma~\ref{lem:quotient-reduced-explicit}, \(R\) is
reduced.  It need not be a field; no factor of \(f\) is selected.

Expand in the polynomial ring \(R[T]\):
\[
             f(u+T)=cT+T^2D(T),\qquad c=f'(u).
\]
The constant term is zero because \(f(u)=0\).  Evaluating
\eqref{eq:squarefree-bezout} at \(u\) gives
\(V(u)c=1\).  Thus \(c\) is a unit with a supplied inverse.
This is the precise substitute for choosing a simple rational root.

For each nonzero row \(A_j\), the polynomial \(A_j(u+T)\) is nonzero:
its highest coefficient in \(T\) is the original nonzero leading scalar
of \(A_j\).  Let
\[
 A_j(u+T)=b_jT^{a_j}+\text{terms of larger \(T\)-degree},
 \qquad b_j\ne0.
\]
The least nonzero coefficient is found by a finite scan.  Define
\(w_j=ma_j+j\).  Distinct indices \(j\in\{0,\ldots,m-1\}\) have
distinct residues modulo \(m\), so the minimum \(w_j\) is attained at
exactly one index \(k\).  Write \(a=a_k\), \(b=b_k\).

\begin{lemma}[The initial coefficient of the local norm]
\label{lem:local-norm-initial}
The coefficient of \(T^{ma+k}\) in the norm after substituting
\(X=u+T\) is
\begin{equation}
 \label{eq:initial-local-norm}
             b^m\bigl((-1)^{m-1}c\bigr)^k,
\end{equation}
and every coefficient of smaller degree is zero.
\end{lemma}

\begin{proof}
Use the basis \(1,Y,\ldots,Y^{m-1}\), and index a matrix column by
\(i\in\{0,\ldots,m-1\}\).  A contribution from \(A_jY^j\) sends
that column to row \(r\equiv i+j\pmod m\), with factor
\(f(u+T)^q\), where \(i+j=r+mq\) and \(q\in\{0,1\}\).

For a determinant permutation \(r=\sigma(i)\), write the selected row
index of \(G\) as \(j_i\), with corresponding carry \(q_i\).
Summing the equalities \(i+j_i=\sigma(i)+mq_i\) gives
\[
                  \sum_i j_i=m\sum_i q_i.
\]
Its determinant term has \(T\)-order at least
\[
 \sum_i a_{j_i}+\sum_iq_i
       =\frac1m\sum_i(ma_{j_i}+j_i)
       =\frac1m\sum_iw_{j_i}.
\]
Here \(1/m\) is a statement about natural integer counts, not division by
an element of \(K\).  Terms containing a zero row coefficient vanish.

Every retained weight is at least \(w_k=ma+k\).  If any selected index
differs from \(k\), the average is strictly larger than \(w_k\), and
its integral value is at least \(w_k+1\).  Therefore only the choice
\(j_i=k\) in every column contributes to degree \(w_k\).
It is the unique permutation \(\sigma(i)=i+k\pmod m\).
The norm of multiplication by \(Y\) is
\((-1)^{m-1}f(u+T)\); equivalently this cyclic permutation has sign
\((-1)^{k(m-1)}\), with total carry \(\sum_iq_i=k\).
Its leading coefficient is exactly \eqref{eq:initial-local-norm}.
The argument is a finite determinant expansion over the commutative ring
\(R\); it does not presume that \(R\) has no zero divisors.
\end{proof}

\begin{proof}[Proof of Theorem~\ref{thm:squarefree-norm-producer}]
Since \(b\ne0\) and \(R\) is reduced, \(b^m\ne0\).
Since \(c\) and \(-1\) are units, the initial coefficient
\eqref{eq:initial-local-norm} is nonzero.  Thus the substituted norm is
nonzero.  Determinants commute with coefficient homomorphisms, so this
substituted polynomial is the image of the original norm under
\(X\mapsto u+T\).  A zero original norm would have zero image, a
contradiction.  This proves nonvanishing.

For the degree estimate use the same determinant-permutation expansion
over \(K[X]\), now with \(e_j=\deg A_j\).  Each nonzero determinant
term has degree at most
\[
 \sum_i e_{j_i}+d\sum_iq_i
       =\frac1m\sum_i(me_{j_i}+dj_i)\leq W.
\]
Cancellation can only decrease degree.  This proves the bound without
requiring uniqueness of the largest global weight.
\end{proof}

\paragraph{Example beyond coprime degree weights and rational local roots.}
Over \(\mathbb Q\), take \(m=2\), \(f=X^4+1\), and
\(G=Y-X^2\).  The two global weights are both \(4\), and
\(\gcd(2,4)=2\).  The polynomial \(f\) has no rational root.
Nevertheless it is squarefree and
\[
 N(G)=\det\begin{pmatrix}-X^2&X^4+1\\1&-X^2\end{pmatrix}=-1.
\]
In the quotient algebra, \(u^4=-1\), so the least local row is the
constant \(-u^2\), whose square is \(-1\).  This example shows why the
quotient-root construction is a real enlargement of the earlier guards.

The zero-divisor qualification is not cosmetic.  For
\(f=X^2-1\), \(m=2\), and \(G=(X-1)+Y\), the least local coefficient
is \(b=u-1\).  It is a nonzero zero divisor in
\(\mathbb Q[u]/(u^2-1)\), but \(b^2=2-2u\ne0\).  Indeed
\(N(G)=2-2X\).  A proof that silently inverted \(b\), or treated
this coefficient algebra as a field, would miss precisely this case.

\paragraph{Necessary qualifications.}
Nonconstancy is essential: the constant squarefree polynomial \(f=1\)
would give a zero quotient algebra and reducible \(Y^2-1\).
Squarefreeness is a sufficient condition, not a necessary one; without
it the present automatic rule is unavailable.  For example,
\(f=X^2\), \(m=2\), and \(G=Y-X\) have zero norm.
The nonvanishing statement concerns nonzero polynomials reduced to
\(Y\)-degree below \(m\); the defining equation itself has zero class.

\subsection{Finite and object-arithmetic obligations}

For a uniform primitive-recursive claim, ``exact field'' must mean a
specified coefficient presentation with primitive-recursive validity,
equality, addition, multiplication, additive inverse, and inversion on
nonzero valid codes, together with their arithmetic proofs.  Decidable
equality alone would not establish these complexity bounds.  Over
\(\mathbb Q\), signed numerators and positive denominators provide the
standard presentation.  A schema over arbitrary fields remains a
metatheorem until such a presentation is supplied.

The finite implementation and its proof have the following obligations.
\begin{enumerate}
\item Compute degrees and leading coefficients by bounded scans of the
      input arrays.  Normalize the divisor for remainder computations,
      retaining the original leading coefficient for \(f(u+T)\).
\item Compute polynomial gcd and Bezout certificates by Euclidean
      division.  Degree decreases bound the number of nonzero remainders;
      all coefficients and intermediate arrays arise from explicitly
      bounded field-operation loops.  Verify the displayed Bezout
      identities rather than invoking a factorization oracle.
\item Represent \(R\) by residue arrays of length \(d\).  Equality is
      equality of canonical remainder arrays.  Addition and multiplication
      reduce modulo the normalized \(f\).  The inverse of \(c\) is the
      specific residue \(V(u)\); inversion of arbitrary nonunits of
      \(R\) is neither available nor needed.
\item Expand \(A_j(u+T)\) and \(f(u+T)\) by finite binomial or Horner
      operations; find the least nonzero coefficient by a bounded scan.
      The weight bound is \(m\max_j\deg A_j+(m-1)\).
\item Compute determinants by the Leibniz formula, or another
      division-free determinant algorithm.  Gaussian elimination with
      inverses of arbitrary quotient-algebra pivots would be invalid.
      The permutation loop has the explicit finite bound \(m!\).
\item Prove the carry-sum identity and the initial-coefficient identity
      by finite sums and bounded permutation bookkeeping.  Reducedness
      uses the explicit certificate of Lemma~\ref{lem:quotient-reduced-explicit};
      its geometric sum has length \(m\) when applied to \(b^m\).
\item Export the nonzero coefficient as a finite witness, then compose
      with the existing coefficient-homomorphism theorem for determinants
      and the degree-bound theorem.  The endpoint is an object-iPRA
      implication from the exact-field and squarefree certificates to
      a nonzero norm coefficient and the claimed degree cap.
\end{enumerate}
All dimensions, exponents, and polynomial lengths may be numerical inputs
to fixed primitive-recursive array programs.  No unbounded factor search,
choice of an irreducible factor, extension-field existence theorem,
analytic limit, or division by \(m\) in the coefficient field occurs.
This audit supplies the full human proof and a direct finite formalization
route; it does not itself claim a compiled Lean object-theory theorem.
